\documentclass[../candidatura.tex]{subfiles}
\graphicspath{{\subfix{../imagenes/}}}
\begin{document}

\section{Objetivos de la investigación}

% en borradores iniciales, la extensión sugerida era de media página

\subsection{Objetivo general}

\begin{enumerate}
    \item Desarrollar una propuesta integral de popusintesíntesis mediante la construcción de un glosario conceptual y metodológico, y el diseño y fabricación de infraestructura electrónica y computacional, con el fin de promover nuevas formas de luthería popular, artesanal y abierta para la creación de sintetizadores sonoros.
\end{enumerate}

\subsection{Objetivos secundarios}

\begin{enumerate}
    \item Delimitar la popusintesíntesis como campo de prácticas y saberes, articulando dimensiones textuales, metafóricas, gestuales \footcite{haudricourt_cultivo_2019}, electrónicas y computacionales en un glosario que haga visibles y transmisibles conocimientos que hoy circulan de forma fragmentada y con altas barreras técnicas.
    \item Situar críticamente la propia práctica artística en diálogo con referentes artísticos, productivos y bibliográficos de la luthería electrónica abierta y artesanal, para reconocer cuáles estrategias han permitido, o impedido, que otras personas construyan sus propios sintetizadores.
    \item Pesquisar cómo las decisiones materiales en electrónica, como el uso de componentes de bajo costo, des soldadura sin plomo, la legibilidad de las placas de circuito impreso y criterios de sustentabilidad y reparabilidad, inciden en quiénes pueden construir, reparar y modificar un sintetizador.
    \item Examinar cómo las decisiones computacionales, como la elección de lenguajes de programación, el diseño de bibliotecas de código y la automatización, abren o cierran el diseño de sintetizadores a personas sin formación técnica especializada.
    \item Analizar cómo la documentación, la docencia y el trabajo público permiten que estos conocimientos sean apropiados, reproducidos y adaptados por otras personas en distintos contextos productivos y culturales.
\end{enumerate}

\end{document}